% 恒星光谱（综述）
% license CCBYSA3
% type Wiki

本文根据 CC-BY-SA 协议转载翻译自维基百科\href{https://en.wikipedia.org/wiki/Stellar_classification}{相关文章}。

在天文学中，恒星分类是根据恒星的光谱特性对恒星进行的分类。来自恒星的电磁辐射通过棱镜或衍射光栅被分解成光谱，呈现出类似彩虹的连续颜色，其间夹杂着光谱线。每一条光谱线都对应一种特定的化学元素或分子，而谱线的强度则反映该元素的丰度。不同光谱线强度的变化主要由恒星光球层的温度决定，尽管在某些情况下也确实存在真实的元素丰度差异。恒星的光谱型是一种简短的代码，主要概括了其电离状态，从而为光球层温度提供一个客观的度量。

目前，大多数恒星采用摩根–基南（Morgan–Keenan，MK）分类系统进行分类，该系统使用字母 O、B、A、F、G、K、M，按从最热（O 型）到最冷（M 型）的顺序排列。每一个字母类别又进一步用一个数字细分，其中 0 表示最热，9 表示最冷（例如 A8、A9、F0、F1 就构成了一个由热到冷的序列）。此外，该序列还扩展出 W、S、C 三个类别，用于那些不符合经典体系的恒星。一些恒星遗迹或质量明显偏离恒星范围的天体也被赋予了字母标记：D 表示白矮星，L、T、Y 表示棕矮星（以及部分系外行星）。

在 MK 系统中，还会在光谱型后加上一个用罗马数字表示的光度级。光度级依据恒星光谱中某些吸收线的宽度来判定，这些线宽随恒星大气密度变化，因此可用于区分巨星与矮星。其中，光度级 0 或 Ia+ 表示特超巨星，I 表示超巨星，II 表示亮巨星，III 表示普通巨星，IV 表示亚巨星，V 表示主序星，sd（或 VI）表示亚矮星，而 D（或 VII）表示白矮星。太阳的完整光谱分类为 G2V，这表明太阳是一颗主序星，其表面温度约为 5800 K。
\subsection{传统颜色描述}
传统的恒星颜色描述只考虑恒星光谱的峰值位置。然而实际上，恒星会在整个电磁波谱范围内辐射。由于所有光谱颜色叠加后呈现为白色，人眼实际看到的恒星颜色要比传统颜色描述所暗示的浅得多。这种“明亮度”的特性说明，将恒星简单地赋予某种光谱颜色可能具有误导性。在排除暗光条件下的颜色对比效应后，在一般观测条件中，并不存在真正的绿色、青色、靛色或紫色恒星。像太阳这样的“黄色”矮星实际上呈现为白色；“红矮星”呈现为较深的黄色或橙色；而“棕矮星”并不会真的呈现棕色，而是假想中在近距离观察时可能显得暗红色，或呈灰色/黑色。
\subsection{现代分类}
现代恒星分类体系被称为摩根–基南（Morgan–Keenan，MK）分类系统。在该系统中，每一颗恒星都会被赋予一个光谱型（源自较早的哈佛光谱分类体系，该体系本身并不包含光度级$^{[1]}$），并再配以一个用罗马数字表示的光度级（如下所述），二者共同构成恒星的完整光谱类型。

其他现代恒星分类体系（例如 UBV 系统）则基于颜色指数，即通过测量两个或多个波段星等之间的差值来进行分类$^{[2]}$。这些数值通常以诸如 U−V、B−V这样的形式给出，表示通过两种标准滤光片（例如紫外、蓝光和可见光）所测得的颜色差异。
\subsubsection{哈佛光谱分类}
哈佛体系是一种一维的恒星分类方案，由天文学家安妮·江普·坎农（Annie Jump Cannon）提出，她对德雷珀（Draper）早期按字母顺序排列的分类体系进行了重新排序与简化（见“历史”部分）。在该体系中，恒星根据其光谱特征被划分为若干组，每一组用一个字母表示，并可辅以数字细分。

主序星的表面温度大致分布在 2000–50 000 K 之间，而演化程度更高的恒星——尤其是新近形成的白矮星——其表面温度则可超过 100 000 K$^{[3]}$。从物理意义上看，这些光谱类别反映的是恒星大气的温度，并通常按由热到冷的顺序排列。
\begin{figure}[ht]
\centering
\includegraphics[width=14.25cm]{./figures/1615d7b930f17d79.png}
\caption{} \label{fig_Stcl_1}
\end{figure}
用于记忆光谱型字母由热到冷排列顺序的传统助记口诀是“Oh, Be A Fine Guy/Girl: Kiss Me!”（“哦，做个好小伙/好姑娘：吻我吧！”）$^{[12]}$。
尽管在天文学课程和相关组织举办的活动中，人们提出过许多替代性的助记口诀，但这一传统口诀至今仍然是最为流行的$^{[13][14]}$。

光谱类型 O 到 M（以及后文讨论的其他更为专门的类型）都会再用阿拉伯数字 0–9进行细分，其中 0 表示该类型中最热的恒星。例如，A0 表示 A 型中最热的恒星，而 A9 则表示最冷的 A 型恒星。允许使用小数；例如，恒星 Mu Normae 的光谱型被定为 O9.7 $^{[15]}$。太阳的光谱型为 G2$^{[16]}$。

尽管哈佛光谱分类实际上反映的是恒星的表面温度或光球温度（更准确地说是其有效温度），但这一点在该体系建立之初并未被完全理解。不过，到第一张赫罗图（Hertzsprung–Russell diagram）于 1914 年前后提出时，人们已经普遍怀疑这一对应关系是成立的$^{[17]}$。在 1920 年代，印度物理学家梅格纳德·萨哈（Meghnad Saha）在物理化学中分子解离理论的基础上，发展出了电离理论并将其推广到原子的电离问题。他首先将该理论应用于太阳色球层，随后又应用于恒星光谱$^{[18]}$。

随后，哈佛天文学家塞西莉亚·佩恩（Cecilia Payne）证明O–B–A–F–G–K–M这一光谱序列本质上就是一个温度序列$^{[19]}$。由于该分类序列的建立早于人们对其“温度本质”的认识，将一条光谱归入某一具体亚型（如 B3 或 A7）在很大程度上依赖于对恒星光谱中吸收特征强度的（带有一定主观性的）估计。因此，这些亚型并不是在数学意义上等间隔划分的。
\subsubsection{摩根–基南（Morgan–Keenan）分类}
耶基斯光谱分类，也称为 MK 分类，或 Morgan–Keenan（亦称 MKK，Morgan–Keenan–Kellman）系统$^{[20][21]}$，是一种由威廉·威尔逊·摩根（William Wilson Morgan）、菲利普·C·基南（Philip C. Keenan）和伊迪丝·凯尔曼（Edith Kellman）于 1943 年在耶基斯天文台提出的恒星光谱分类体系$^{[22]}$。这是一个二维分类体系（温度与光度），其依据是对恒星温度和表面重力敏感的光谱线；而表面重力又与恒星的光度密切相关（相比之下，哈佛分类仅基于表面温度）。在 1953 年对标准星表和分类标准进行若干修订之后，该体系被正式命名为摩根–基南分类（MK）$^{[23]}$，并一直沿用至今。

表面重力更高、密度更大的恒星，其光谱线会表现出更显著的压力展宽。由于巨星的半径远大于质量相近的矮星，其表面重力（以及由此产生的压力）要低得多。因此，光谱中的差异可以被解释为光度效应，从而仅通过分析光谱本身就能够确定恒星的光度级。

目前区分出了若干不同的光度级，具体列于下表中$^{[24]}$。
\begin{figure}[ht]
\centering
\includegraphics[width=14.25cm]{./figures/73151fcf04a39c9e.png}
\caption{} \label{fig_Stcl_2}
\end{figure}
允许存在边缘情况；例如，一颗恒星可能既可被归为超巨星也可被归为亮巨星，或者介于亚巨星与主序星之间。在这些情况下，会在两个光度级之间使用两种特殊符号：
\begin{itemize}
\item 斜杠（/）：表示恒星可能属于其中任意一个光度级。
\item 连字符（-）：表示恒星处在两个光度级之间的过渡状态。
\end{itemize}
例如，被标记为 A3-4III/IV 的恒星，表示其光谱型介于 A3 与 A4 之间，同时其光度级介于巨星（III）与亚巨星（IV）之间，或可能属于其中之一。

此外，还曾使用亚矮星（sub-dwarf）光度级：

VI 表示亚矮星（其光度略低于主序星）。

名义上的光度级 VII（有时甚至更高的罗马数字）如今已很少用于白矮星或“热亚矮星”的分类，因为主序星和巨星所采用的温度字母体系已不再适用于白矮星。

在某些情况下，字母 a 和 b 也会用于非超巨星的光度级中；例如，一颗光度略低于典型巨星的恒星可被标为 IIIb，而 IIIa 则表示其光度略高于典型巨星$^{[34]}$。

一小部分 V 型极端恒星在 He II λ4686 吸收谱线处表现出很强的吸收特征，因此被赋予 Vz 的分类标记。一个典型例子是恒星 HD 93129 B$^{[35]}$。
\subsubsection{光谱异常}
还可以在光谱型之后附加小写字母形式的补充标记，用以指出光谱中的某些特殊或异常特征$^{[36]}$。
\begin{table}[ht]
\centering
\caption{请输入表格标题}\label{tab_Stcl1}
\begin{tabular}{|c|c|}
\hline
\textbf{代码} & \textbf{恒星的光谱异常} \\
\hline
: & 不确定的光谱值$^{[24]}$\\
\hline
... & 存在未描述的光谱异常\\
\hline
! & 特殊异常\\
\hline
comp & 复合光谱[37]\\
\hline
e & 存在发射线[37]\\
\hline
[e] & 存在“禁戒”发射线\\
\hline
er & “反向”发射线中心比边缘弱\\
\hline
eq & 具有P Cygni轮廓的发射线\\
\hline
f & N III 和 He II 发射$^{[24]}$\\
\hline
f & NIV 4058Å线比NIII 4634Å、4640Å和4642Å线更强$^{[38]}$\\
\hline
f+ & 除了N III线外，还发射Si IV 4089Å和4116Å线$^{[38]}$\\
\hline
f? & C III 4647–4650–4652Å发射线的强度与N III线相当$^{[39]}$\\
\hline
(f) & N III发射，He II缺失或吸收较弱\\
\hline
(f+) & $^{[40]}$\\
\hline
((f)) & 显示强烈的He II吸收，伴随微弱的N III发射$^{[41]}$\\
\hline
((f*)) & $^{[40]}$\\
\hline
h & WR恒星以氢发射线为特征。$^{[42]}$\\
\hline
ha & WR恒星的氢线在吸收和发射中均可见。$^{[42]}$\\
\hline
He wk & 弱氦线\\
\hline
k & 带有星际吸收特征的光谱\\
\hline
m & 增强的金属特征$^{[37]}$\\
\hline
n & 因自旋引起的宽泛（“模糊”）吸收$^{[37]}$\\
\hline
nn & 非常宽泛的吸收特征$^{[24]}$\\
\hline
neb & 星云的光谱混入其中$^{[37]}$\\
\hline
p & 未指定的特殊性，奇特恒星.\\
\hline
pq & 奇特光谱，类似于新星的光谱\\
\hline
q & P型仙后座轮廓\\
\hline
s & 窄（“锐利”）吸收线$^{[37]}$\\
\hline
ss & 极窄线\\
\hline
sh & 壳层星特征$^{[37]}$\\
\hline
变量 & 可变光谱特征$^{[37]}$（有时简称为“v”）\\
\hline
波长 & 弱线$^{[37]}$（也称“w”和“wk”）\\
\hline
元素符号 & 指定元素的异常强烈的光谱线$^{[37]}$\\
\hline
z & 表明在468.6 纳米处存在异常强烈的电离氦线$^{[35]}$\\
\hline
\end{tabular}
\end{table}
例如，59 Cygni 被列为光谱型 B1.5Vnne$^{[43]}$，这表示其光谱的总体分类为 B1.5V，同时还具有非常宽的吸收线以及某些发射线。
\subsection{历史}
哈佛光谱分类中字母顺序之所以显得有些“反常”，其原因在于历史演变：该体系起源于更早的塞基（Secchi）分类，并随着人们对恒星物理认识的不断加深而逐步修订和重组。
\subsubsection{塞基分类}
在 19 世纪 60—70 年代，恒星光谱学的先驱安杰洛·塞基（Angelo Secchi） 创建了塞基光谱分类体系，用于对观测到的恒星光谱进行分类。到 1866 年，他已经提出了三类恒星光谱，如下表所示$^{[44][45][46]}$。

在 19 世纪末（1890 年代后期），这一分类体系逐渐被哈佛光谱分类所取代；后者正是本文其余部分所讨论的现代恒星光谱分类基础$^{[47][48][49]}$。
\begin{table}[ht]
\centering
\caption{请输入表格标题}\label{tab_Stcl2}
\begin{tabular}{|c|c|}
\hline
\textbf{等级编号} & \textbf{塞奇等级描述}\\
\hline
塞奇等级I & 白蓝相间的恒星，具有宽阔而强烈的氢线，例如织女星和牛郎星。这包括现代A型和早期F型。\\
\hline
塞奇等级I（猎户座亚型）& $<b$一种属于Secchi I类的亚型，其光谱中以窄线取代了宽带，例如参宿四和大角星。按现代术语来说，这对应于早期B型恒星\\
\hline
塞奇等级II & 黄色恒星——氢线较弱，但金属谱线明显，例如太阳, 大角星和五车二。这包括现代的G型和K型恒星，以及晚期F型恒星。\\
\hline
塞奇等级III & Orange to red stars with complex band spectra, such as Betelgeuse and Antares.
这对应于现代的M类。\\
\hline
塞奇四等 & In 1868, he discovered carbon stars, which he put into a distinct group:$^{[50]}$带有显著碳谱带和谱线的红色恒星，对应于现代的C类和S类。\\
\hline
塞奇五等 & 1877年，他增设了第五等：$^{[51]}$发射线星，例如仙后座γ星和谢利阿克星, 属于现代的Be型。1891年，爱德华·查尔斯·皮克林提出，V类应对应于现代的O型（当时还包括沃尔夫-拉叶星）。以及行星状星云中的恒星。\\
\hline
\end{tabular}
\end{table}
用于塞基（Secchi）光谱分类的罗马数字，不应与耶基斯（Yerkes）光度分类以及拟议的中子星分类中所使用的、彼此毫不相关的罗马数字混淆。
\subsubsection{德雷珀体系}
\begin{figure}[ht]
\centering
\includegraphics[width=14.25cm]{./figures/a92bc94a88248741.png}
\caption{} \label{fig_Stcl_3}
\end{figure}
在丈夫去世后，玛丽·安娜·德雷珀（Mary Anna Draper）开始资助哈佛底片库（Harvard Plate Stacks）的建立，以及在哈佛学院天文台对这些天文底片的研究。天文台台长爱德华·C·皮克林（Edward C. Pickering）随后雇用了多位开创性的女性天文学家，她们后来被统称为“哈佛计算员”（Harvard Computers）。尽管最初设想她们会研究多个不同的天文学课题，但这项工作的一个早期重要成果，是《亨利·德雷珀恒星光谱纪念目录》（The Henry Draper Memorial Catalogue of Stellar Spectra）的第一版，该目录于 1890 年首次出版。

威廉敏娜·弗莱明（Williamina Fleming）完成了该目录第一版中大部分光谱的分类工作，她被认为共分类了 一万余颗恒星，并发现了 10 颗新星以及 200 多颗变星。在哈佛计算员，尤其是威廉敏娜·弗莱明的帮助下，亨利·德雷珀目录的最初版本被设计出来，用以取代由安杰洛·塞基（Angelo Secchi）建立的罗马数字光谱分类体系。

该目录采用了一种新的分类方案：将此前使用的塞基分类（I 至 V 类）进一步细分为更具体的类别，并用 A 到 P 的字母表示；此外，还使用字母 Q 来标记那些无法归入其他任何类别的恒星。$^{[53][54]}$弗莱明与皮克林合作，根据氢谱线强度的不同区分出了 17 个光谱类别。氢谱线强度的变化会导致恒星辐射波长的差异，从而造成颜色外观的变化。在这一体系中，A 类光谱表现出最强的氢吸收线，而 O 类光谱几乎不显示可见谱线。随着字母顺序向后推进，氢吸收线的强度逐渐减弱。后来，安妮·跳·坎农（Annie Jump Cannon）和安东尼娅·莫里（Antonia Maury）对该体系进行了修订，最终形成了哈佛光谱分类体系。$^{[56][58]}$
\subsubsection{旧哈佛体系（1897）}
1897 年，哈佛的另一位天文学家安东尼娅·莫里将塞基 I 类中的“猎户型”亚类置于塞基 I 类其余部分之前，从而在分类顺序上使现代的 B 型位于现代的 A 型之前。她是第一位这样做的人，尽管她并未使用字母标记的光谱类型，而是采用了一个由 I–XXII 共 22 个编号组成的分类序列。$^{[59][60]}$
\begin{table}[ht]
\centering
\caption{1897年哈佛系统的概要$^{[61]}$}\label{tab_Stcl3}
\begin{tabular}{|c|c|}
\hline
\textbf{分组} & \textbf{摘要}\\
\hline
I−V & 包括了从I组到V组氢吸收线强度逐渐增强的‘猎户座型’恒星\\
\hline
VI & 充当了‘猎户座型’与Secchi I型之间的中间类型\\
\hline
VII−XI & 是塞基1型恒星，其氢吸收线强度从第VII组至第XI组逐渐减弱\\
\hline
XIII−XVI & 包括塞基2型恒星，其氢吸收线强度逐渐减弱，而太阳型金属线则逐渐增强\\
\hline
XVII−XX & $<b$包括光谱线逐渐增多的塞奇3型恒星\\
\hline
XXI & 包括第4型天狼星\\
\hline
二十二 & 包括沃尔夫-拉叶星\\
\hline
\end{tabular}
\end{table}
由于最初的 22 个罗马数字分组仍不足以涵盖恒星光谱中出现的全部变化，人们又引入了三种附加划分，以进一步细化差异：在原有编号后添加小写字母，用来区分谱线外观的相对特征；这些谱线被定义为：
\begin{itemize}
\item $a$：平均宽度
\item $b$：模糊
\item $c$：锐利
\end{itemize}
安东尼娅·莫里于 1897 年发表了她自己的恒星分类目录，题为
《用 11 英寸德雷珀望远镜拍摄的亮星光谱：亨利·德雷珀纪念项目的一部分》（Spectra of Bright Stars Photographed with the 11 inch Draper Telescope as Part of the Henry Draper Memorial）。
该目录收录了 4,800 张照片以及莫里本人对 681 颗北天亮星的分析。这是首次由女性署名的天文台正式出版物**。
\subsubsection{现行哈佛体系（1912）}
1901 年，安妮·跳·坎农重新采用了字母分类法，但她舍弃了除 O、B、A、F、G、K、M、N 之外的所有字母，并按这一顺序使用；此外，她还保留了 P（用于行星状星云）和 Q（用于某些异常光谱）。她还引入了诸如 B5A（介于 B 型与 A 型之间的恒星）、F2G（从 F 型向 G 型过渡五分之一）的混合标记方式，依此类推。

最终，到 1912 年，坎农将 B、A、B5A、F2G 等类型统一改写为 B0、A0、B5、F2 等形式。这一体系基本上就是现代哈佛光谱分类系统。该系统建立在对摄影底片上恒星光谱的系统分析之上，这种方法能够将恒星发出的光转换为可读的光谱。
\subsubsection{威尔逊山分类}
一种称为威尔逊山体系的光度分类法曾被用于区分不同光度的恒星。即便在今天，这种标记方式仍偶尔可在现代光谱中见到。其主要符号包括：
\begin{itemize}
\item sd：次矮星（subdwarf）
\item d：矮星（dwarf）
\item sg：次巨星（subgiant）
\item g：巨星（giant）
\item c：超巨星（supergiant）
\end{itemize}
\subsection{光谱型}
“光谱型”（Spectral type）一词在此指恒星光谱分类；若指小行星的光谱分类，则另见“小行星光谱类型”。恒星分类体系本质上是一种分类学（taxonomy）体系，以模式恒星为基础，类似于生物学中对物种的分类：每一个类别及其子类别都由一颗或多颗标准恒星来定义，并配有相应的判别特征描述。
\subsubsection{“早型（early）”与“晚型（late）”的命名法}
恒星常被称为早型或晚型。其中，“早型”是更热的同义词，而“晚型”是更冷的同义词。

根据具体语境，“早”和“晚”既可以是绝对概念，也可以是相对概念。作为绝对概念时，“早型”通常指 O 型或 B 型恒星，有时也包括 A 型恒星；作为相对概念时，“早型”则指在同一光谱型内部温度较高的恒星，例如“早期 K 型”可能指 K0、K1、K2 或 K3。

“晚型”的用法与此类似。不加限定地使用“晚型”时，通常指 K 型和 M 型恒星；但在相对意义上，它也可用于表示在同一光谱型中相对较冷的恒星，例如“晚期 G 型”指 G7、G8 和 G9。

在这种相对用法中，“早型”对应于光谱型字母后较小的阿拉伯数字，“晚型”则对应于较大的数字。

这种看似晦涩的术语源自 19 世纪末的一种恒星演化模型残余。该模型认为恒星的能量来自引力收缩，即开尔文–亥姆霍兹机制；而如今已知，这一机制并不适用于主序星。按照该假设，恒星会在生命初期是非常炽热的“早型星”，随后逐渐冷却，演化为“晚型星”。然而，这一机制给出的太阳年龄远小于地质记录所显示的真实年龄，并在发现恒星能量来源于核聚变之后被彻底否定。尽管如此，“早型”和“晚型”这两个术语仍被沿用下来，尽管它们所依托的理论早已被淘汰。
\subsubsection{O 型（Class O）}
\begin{figure}[ht]
\centering
\includegraphics[width=8cm]{./figures/2338d2de317ddeae.png}
\caption{一颗假想的 O5V 型恒星的光谱*} \label{fig_Stcl_4}
\end{figure}
O 型恒星极其炽热且光度极高，其辐射能量的大部分位于紫外波段。它们是所有主序星中最为稀少的一类。在太阳邻域的主序星中，大约每 3,000,000 颗才有 1 颗（约 0.00003\%）属于 O 型星 $^{[c][11]}$。一些质量最大的恒星就位于这一光谱型内。O 型恒星周围的环境往往十分复杂，这也使得其光谱测量变得困难。

早期对 O 型光谱的定义基于 He II $\lambda4541$与 He I $\lambda4471$ 谱线强度之比，其中 $\lambda$ 表示辐射波长。光谱型 O7 被定义为这两条谱线强度相等的位置，而在更早的亚型中，He I 谱线会逐渐减弱。按最初定义，O3 型对应于该谱线完全消失的情形，不过在现代技术条件下仍可极其微弱地观测到它。因此，现代定义改为采用氮谱线 N IV $\lambda4058$ 与 N III $\lambda\lambda4634$–$40$–$42$的强度比 $^{[74]}$。

O 型恒星的光谱以 He II 吸收线（有时也伴随发射线）为主，同时具有显著的电离元素谱线（如 Si IV、O III、N III、C III）以及中性氦谱线，这些氦谱线在 O5 到 O9 之间逐渐增强；还可见明显的氢巴耳末系谱线，但其强度不及更晚型恒星。质量更高的 O 型恒星由于其恒星风速度极高（可达 2,000 km/s），难以保留厚重的大气层。正因为质量巨大，O 型恒星的核心温度极高，氢燃料消耗得非常迅速，因此它们往往是最早离开主序带的恒星。

当 MKK 分类方案于 1943 年首次提出时，O 型仅包含 O5–O9.5 这些亚型 $^{[75]}$。随后，该方案在 1971 年扩展到 O9.7 $^{[76]}$，在 1978 年扩展到 O4$^{[77]}$，并在之后又引入了 O2、O3 以及 O3.5 等新的亚型 $^{[78]}$。

光谱标准星示例 $^{[72]}$：
\begin{itemize}
\item O7V —— 麒麟座 S（S Monocerotis）
\item O9V —— 蝎虎座 10（10 Lacertae）
\end{itemize}
\subsubsection{B 型恒星（Class B）}
\begin{figure}[ht]
\centering
\includegraphics[width=8cm]{./figures/a4b2537f0543257b.png}
\caption{假想的 B3V 型恒星的光谱} \label{fig_Stcl_5}
\end{figure}
B 型恒星非常明亮，呈蓝色。它们的光谱中包含中性氦谱线，在 B2 亚型中最为显著，同时还具有中等强度的氢谱线。由于 O 型和 B 型恒星能量极高，它们的寿命相对较短。因此，在其一生中发生显著动力学相互作用的概率很低，除了一些“逃逸星”（runaway stars）之外，它们通常不会远离自身形成的区域。

最初，O 型向 B 型的过渡被定义为 He II λ4541 谱线消失的界限。然而，借助现代观测设备，这条谱线在早期 B 型恒星中仍然可以观测到。如今，对于主序星而言，B 型恒星的定义改为依据 He I 紫端谱线的强度，其最大强度对应于 B2 型。对于超巨星，则使用硅谱线进行判定：Si IV λ4089 与 Si III λ4552 的谱线表明其为早期 B 型；在中期 B 型中，后者相对于 Si II λλ4128–30 的强度是关键判据；而在晚期 B 型中，则以 Mg II λ4481 相对于 He I λ4471 的强度作为区分标准。${}^{[74]}$

这些恒星往往分布在其起源的 OB 星协中，而 OB 星协通常与巨分子云相关联。猎户座 OB1 星协占据了银河系一条旋臂的相当大一部分，并包含猎户座中许多较亮的恒星。在太阳邻域的主序星中，大约每 800 颗恒星中就有 1 颗（约 0.125\%）是 B 型主序星。${}^{[c][11]}$B 型恒星相对少见，其中距离我们最近的是轩辕十四（Regulus），约 80 光年之外。${}^{[79]}$

一些质量很大但并非超巨星的恒星被称为 Be 星。这类恒星的光谱中可观测到一条或多条巴耳末系发射线，其中由恒星向外投射的与氢相关的电磁辐射系列尤为引人关注。一般认为，Be 星具有异常强烈的恒星风、较高的表面温度，并且由于其以异常快的速度自转，会发生显著的恒星质量流失。${}^{[80]}$

另一类被称为 B[e] 星（或因排版原因写作 B(e) 星）的天体，具有显著的中性或低电离态发射线。这些谱线被认为涉及禁戒跃迁机制，即发生了在现有量子力学理解下通常不被允许的过程。

典型光谱标准星：${}^{[72]}$
\begin{itemize}
\item B0V —— 猎户座 υ 星（Upsilon Orionis）
\item B0Ia —— 参宿二（Alnilam）
\item B2Ia —— 猎户座 χ² 星（Chi² Orionis）
\item B2Ib —— 仙王座 9 星（9 Cephei）
\item B3V —— 北斗七星之七·摇光（Alkaid）
\item B3V —— Haedus
\item B3Ia —— 大犬座 ο² 星（Omicron² Canis Majoris）
\item B5Ia —— 船尾座 ε 星（Aludra）
\item B8Ia —— 参宿七（Rigel）
\end{itemize}
\subsubsection{A 型恒星（Class A）}
\begin{figure}[ht]
\centering
\includegraphics[width=8cm]{./figures/a2e592b0bb3c9ac1.png}
\caption{一颗假想的 A5V 型恒星的光谱} \label{fig_Stcl_6}
\end{figure}
A 型恒星是在肉眼可见恒星中较为常见的一类，颜色呈白色或蓝白色。它们具有很强的氢谱线，在 A0 型达到最强；同时还具有电离金属谱线（如 Fe II、Mg II、Si II），其强度在 A5 型达到最大。到这一阶段，Ca II 谱线的增强也尤为明显。在太阳邻域的主序星中，大约每 160 颗恒星中就有 1 颗（约 0.625\%）是 A 型恒星，${}^{[c][11]}$其中在 15 秒差距范围内就包含 9 颗 A 型恒星。$^{[81]}$

典型光谱标准星：${}^{[72]}$
\begin{itemize}
\item A0Van —— 天仓五（Phecda）
\item A0Va —— 织女星（Vega）
\item A0Ib —— 狮子座 η 星（Eta Leonis）
\item A0Ia —— HD 21389
\item A1V —— 天狼星 A（Sirius A）
\item A2Ia —— 天津四（Deneb）
\item A3Va —— 南鱼座 α 星（Fomalhaut）
\end{itemize}
\subsubsection{F 型恒星（Class F）}
\begin{figure}[ht]
\centering
\includegraphics[width=8cm]{./figures/88e11dd026c5811b.png}
\caption{一颗假想的 F5V 型恒星的光谱} \label{fig_Stcl_7}
\end{figure}
F 型恒星的光谱中，钙的 Ca II H 与 K 线逐渐增强；到晚期 F 型时，中性金属线（如 Fe I、Cr I）开始超过电离金属线。其光谱特征表现为氢线较弱、电离金属线较弱。它们的颜色呈白色。在太阳邻域的主序星中，大约 1/33（3.03\%）属于 F 型恒星$^{[c][11]}$，其中在 20 光年范围内仅有一颗，即南河三 A（Procyon A）$^{[82]}$。

\textbf{典型光谱标准星：}$^{[72][83][84][85][86]}$
\begin{itemize}
\item F0IIIa —— Adhafera
\item F0Ib —— Arneb
\item F1V —— 37 Ursae Majoris（大熊座 37）
\item F2V —— 78 Ursae Majoris（大熊座 78）
\item F7V —— Iota Piscium（双鱼座 ι）
\item F9V —— Zavijava
\item F9V —— HD 10647
\end{itemize}
\subsubsection{G 型}
\begin{figure}[ht]
\centering
\includegraphics[width=8cm]{./figures/d758de7d2c3b6f51.png}
\caption{假想的 G5V 型恒星光谱} \label{fig_Stcl_8}
\end{figure}
G 型恒星（包括太阳在内）具有显著的 Ca II 的 H、K 吸收线，在 G2 亚型时最为明显。它们的氢吸收线比 F 型恒星更弱，但除电离金属线外，还出现中性金属线。在 CN 分子的 G 带处存在一个显著的谱带峰值。G 型主序星约占太阳邻域主序星的 7.5\%，也就是大约每 13 颗主序星中就有 1 颗是 G 型恒星。在 10 pc 范围内共有 21 颗 G 型恒星。

G 型还包含所谓的“黄色演化空隙”（Yellow Evolutionary Void）。超巨星在演化过程中常在 O 或 B 型（蓝色）与 K 或 M 型（红色）之间摆动，但它们在不稳定的黄色超巨星阶段停留的时间通常很短。

\textbf{光谱标准示例：}$^{[72]}$
\begin{itemize}
\item G0V —— Chara
\item G0IV —— Muphrid
\item G0Ib —— Sadalsuud
\item G2V —— 太阳
\item G5V —— Kappa¹ Ceti
\item G5IV —— Mu Herculis
\item G5Ib —— 9 Pegasi
\item G8V —— 61 Ursae Majoris
\item G8IV —— Alshain
\item G8IIIa —— Kappa Geminorum
\item G8IIIab —— Vindemiatrix
\item G8Ib —— Mebsuta
\end{itemize}
\subsubsection{K 型}
\begin{figure}[ht]
\centering
\includegraphics[width=8cm]{./figures/2bdd33a2f5bc7e0e.png}
\caption{假想的 K5V 型恒星光谱} \label{fig_Stcl_9}
\end{figure}
K 型恒星是呈橙色的恒星，温度略低于太阳。它们约占太阳邻域主序星的 12\%$^{[c][11]}$。此外还存在 K 型巨星，其范围从像 RW Cephei 这样的超巨星，到巨星和超巨星（如大角星 Arcturus）；而像半人马座 $\alpha B$这样的橙矮星则属于主序星。

K 型恒星的氢吸收线极其微弱，甚至可能完全不存在，其光谱主要由中性金属线（如 Mn I、Fe I、Si I）构成。在晚 K 型恒星中，会开始出现氧化钛（TiO）分子带。主流理论（基于较低的有害辐射水平以及恒星寿命较长）因此认为：若行星上的生命形式与地球生命直接类比，这类恒星由于其宜居带较宽、且相较于拥有最宽宜居带的恒星其高能有害辐射阶段显著更少，可能是孕育高度演化生命的最优候选恒星之一$^{[88][89]}$。

\textbf{光谱标准示例：}$^{[72]}$
\begin{itemize}
\item K0V —— Alsafi
\item K0III —— Pollux
\item K0III —— Aljanah
\item K2V —— Ran
\item K2III —— 蛇夫座 κ（Kappa Ophiuchi）
\item K3III —— 牧夫座 ρ（Rho Boötis）
\item K5V —— 天鹅座 61 A（61 Cygni A）
\item K5III —— Eltanin
\end{itemize}
\subsubsection{M 型}
\begin{figure}[ht]
\centering
\includegraphics[width=8cm]{./figures/cdd928e0cdd68726.png}
\caption{假想的 M5V 型恒星光谱} \label{fig_Stcl_10}
\end{figure}
M 型恒星（Class M）是目前为止数量最多的一类恒星。在太阳邻域的主序星中，大约 76\%属于 M 型恒星 $^{[c][f][11]}$。然而，M 型主序星（即红矮星）的光度极低，在一般情况下没有一颗亮到可以用肉眼直接观测，除非在极其特殊的条件下。已知最亮的 M 型主序星是拉卡伊 8760（Lacaille 8760），光谱型为 M0V，视星等为 6.7（而在良好观测条件下，通常认为肉眼可见的极限星等约为 6.5），因此几乎不可能再发现更亮的同类恒星。

尽管大多数 M 型恒星是红矮星，但银河系中一些体积最大的已知超巨星同样属于 M 型，例如大犬座 VY（VY Canis Majoris）、仙王座 VV（VV Cephei）、心宿二（Antares）和参宿四（Betelgeuse）。此外，一些体积较大、温度较高的棕矮星也属于晚期 M 型，通常位于 M6.5 到 M9.5 的范围内。

M 型恒星的光谱包含氧化物分子的谱线（在可见光波段中尤以 TiO 为主）以及各种中性金属的谱线，但通常不出现氢的吸收线。在 M 型恒星中，TiO 分子带往往非常强烈，到了大约 M5 型时，几乎主导了整个可见光光谱。到更晚的 M 型阶段，还会出现二价钒氧化物（VO）的分子带。

\textbf{典型光谱标准星：}$^{[72]}$
\begin{itemize}
\item M3V —— 格利泽 581（Gliese 581）
\item M0IIIa —— 米拉赫（Mirach）
\item M2III —— 飞马座 χ（Chi Pegasi）
\item M1–M2Ia–Iab —— 参宿四（Betelgeuse）
\item M2Ia —— 仙王座 μ（Mu Cephei，“赫歇尔的石榴星”）
\end{itemize}
\subsection{扩展光谱类型}
随着新类型恒星的不断发现，天文学中已经引入了若干新的光谱类型 $^{[90]}$。
\subsubsection{高温蓝色发射星类别}
\begin{figure}[ht]
\centering
\includegraphics[width=8cm]{./figures/d8662d48503162d5.png}
\caption{UGC 5797：一座发射线星系，其中正在形成质量巨大、明亮的蓝色恒星$^{[91]}$} \label{fig_Stcl_11}
\end{figure}
一些极高温、呈蓝色的恒星的光谱中，会出现显著的发射线，这些发射线主要来自碳或氮，有时也来自氧。

\textbf{WR 类（或 W 类）：沃尔夫–拉叶星（Wolf–Rayet）}
\begin{figure}[ht]
\centering
\includegraphics[width=6cm]{./figures/4033fec8a3a32df2.png}
\caption{哈勃空间望远镜拍摄的M1-67 星云及其中心的沃尔夫–瑞叶星 WR 124的图像} \label{fig_Stcl_12}
\end{figure}
曾一度被归入 O 型恒星的W 型（或 WR 型）沃尔夫–瑞叶星$^{[92]}$，其显著特征是光谱中缺乏氢谱线。相反，它们的光谱主要由高度电离的氦、氮、碳，有时还有氧的宽发射线所主导。人们认为，这类恒星大多是演化晚期的超巨星，其外层氢包层已被强烈的恒星风剥离，从而直接暴露出炽热的氦壳层。
WR 型又可根据光谱中氮与碳发射线的相对强度（也反映其外层成分）进一步细分$^{[42]}$。

WR 光谱的分类如下$^{[93][94]}$：

\begin{itemize}
\item WN$^{[42]}$ —— 光谱由 N III–V 和 He I–II 谱线主导

\item WNE（WN2–WN5，部分 WN6）：较热，称为“早型”
\item WNL（WN7–WN9，部分 WN6）：较冷，称为“晚型”
\item 有时使用扩展的 WN10 与 WN11 类型，用于 Ofpe/WN9 恒星$^{[42]}$
\item h标记（如 WN9h）：表示存在氢发射线 ha 标记（如 
WN6ha）：表示同时存在氢的发射与吸收
\item WN/C —— 兼具 WN 特征并出现强 C IV 谱线，介于 WN 与 WC 之间$^{[42]}$
\item WC$^{[42]}$ —— 光谱中具有强烈的 C II–IV 谱线
\item WCE（WC4–WC6）：较热，“早型”
\item WCL（WC7–WC9）：较冷，“晚型”
\item WO（WO1–WO4）—— 具有强 O VI 谱线，极为罕见，可视为 WCE 类向极端高温（可达 200 kK 甚至更高）的延伸
\end{itemize}
尽管大多数行星状星云的中心星（CSPNe）呈现 O 型光谱$^{[95]}$，但约 10\%的此类中心星缺乏氢，并表现出 WR 型光谱$^{[96]}$。这些恒星属于低质量恒星，为与高质量的沃尔夫–瑞叶星区分，其光谱类型用方括号标注，例如 [WC]。其中大多数呈现[WC]光谱，少数为[WO]，而[WN]则极为罕见。

\textbf{斜线星}

斜线星是光谱中同时呈现 WN 型特征谱线的 O 型恒星。名称“斜线（slash）”源于其印刷的光谱类型中包含一个斜杠（例如 Of/WNL）$^{[74]}$。

在这类光谱中还发现了一个次级族群：温度更低、处于“中间态”的恒星，记作 Ofpe/WN9$^{[74]}$。这些恒星也曾被称为 WN10 或 WN11，但随着人们认识到它们与其他沃尔夫–瑞叶星在演化路径上的差异，这种称谓已不再常用。近年来，对更为罕见恒星的发现进一步将斜线星的范围扩展至 O2–3.5If/WN5–7，其温度甚至高于最初定义的斜线星 $^{[97]}$。

\textbf{磁性 O 型恒星}

这是一类具有强磁场的 O 型恒星，其指定标记为 Of?p $^{[74]}$。
\subsubsection{冷红矮星与褐矮星类别}
新的光谱类型 L、T 和 Y 被引入，用于对冷恒星的红外光谱进行分类。这一体系既包括红矮星，也包括在可见光波段极其暗弱的褐矮星 $^{[98]}$。

褐矮星是不发生氢核聚变的天体，它们会随着年龄增长而逐渐冷却，因此会**向更晚的光谱类型演化。褐矮星在形成初期通常呈现 M 型光谱，随后会依次冷却并经历 L、T、Y 光谱类型，而且质量越小，冷却速度越快；质量最高的褐矮星在宇宙当前年龄内甚至还未冷却到 Y 型，乃至 T 型。正因如此，对于某些质量与年龄组合的 L–T–Y 类型，其有效温度与光度之间会出现不可区分的重叠，因此无法为这些类型给出明确的温度或光度取值$^{[10]}$。
\subsubsection{L 型}
\begin{figure}[ht]
\centering
\includegraphics[width=6cm]{./figures/a812750df0497cb6.png}
\caption{} \label{fig_Stcl_13}
\end{figure}
L 型矮星之所以被如此命名，是因为它们**比 M 型恒星更冷**，而字母 **L** 在字母顺序上是**最接近 M 的尚未使用字母**。这类天体中，有一部分质量足够大，能够维持**氢核聚变**，因此严格来说仍属于恒星；但**绝大多数质量低于恒星下限**，因而属于**褐矮星**。它们在可见光下呈现出**非常深的红色**，而在**红外波段最为明亮**。其大气温度足够低，使得**金属氢化物**以及**碱金属**的光谱线在其光谱中尤为显著 $^{[99][100][101]}$。

由于巨星具有**较低的表面引力**，含有 **TiO 和 VO** 的凝结物无法形成。因此，**除矮星之外的 L 型恒星**在**孤立环境中不可能形成**。然而，通过**恒星碰撞**的方式，理论上仍有可能形成 **L 型超巨星**；一个例子是 **V838 Monocerotis**，它在发生**高光度红新星爆发**的高峰期曾表现出类似特征。
\subsubsection{T 型}
\begin{figure}[ht]
\centering
\includegraphics[width=6cm]{./figures/b57b3066df5ac4fb.png}
\caption{} \label{fig_Stcl_14}
\end{figure}